% Independent review fragment. No documentclass and no changes to the main draft.
\subsection{The published Chen input and the smaller removal coefficient}
\label{sec:external-published-chen-lower-k-independent-transfer}

\paragraph{Citation and normalization.}
Runbo Li, \emph{On Chen's theorem, Goldbach's conjecture and almost prime
twins II}, Mathematical Reports 28(78) (2026), 1--2, 39--61,
\href{https://imar.ro/journals/Mathematical_Reports/Pdfs/2026/1_2/3.pdf}
{publisher PDF}, states Theorem~1.1 on p.~40 without a conjectural
hypothesis. Its definitions (1) and (3), on p.~39, give exactly
\[
 D_{1,2}(N)\ge\frac{1233}{625}\mathcal H_N,\qquad
 \mathcal H_N=\mathcal C(N)\frac{N}{\log^2N},\qquad
 \mathcal C(N)=\prod_{\ell>2}\left(1-\frac1{(\ell-1)^2}\right)
 \prod_{\substack{\ell\mid N\\\ell>2}}\frac{\ell-1}{\ell-2},
\]
for every sufficiently large even integer. Prime factors are counted
with multiplicity. There is no extra factor two to import from a twin-prime
normalization. For even $N>2$, its $p\le N$ and our $p<N$ are identical;
the possible convention at $N-p=1$ changes the subsequent subtraction
by at most one. This paragraph checks the published statement, not its proof.

\paragraph{The exact additional theorem required from this project.}
Write
\[
 B_b(N)=\#\{p<N:p\text{ prime},\ N-p=rq,\ q^b<r\le q,
                 \ r,q\text{ prime}\}.
\]
The necessary original conclusion is, for every fixed $1/2<b<1$,
\begin{equation}\label{eq:external-lower-k-required-upper}
 B_b(N)\le\left(\frac{67619}{10000}\log(1/b)+o(1)\right)
               \mathcal H_N
 \quad(N\longrightarrow\infty,\ 2\mid N).
\end{equation}
This is a statement about actual prime counts uniformly along even $N$.
It is not the numerical statement that a displayed integral is below
$67619/10000$.

The finite-product theorem under review specializes to this domain
without arbitrary arithmetic weights. For fixed small $\kappa>0$,
discard $rq<N^{1-\kappa}$, costing $O(N^{1-\kappa})$ prime counts.
On the retained region,
\[
 r,q\ge N^{(1-\kappa)b/(1+b)},
\]
so a fixed positive factor cutoff is available. The inequalities
$rq\le N$, $r\le q$ and $r>q^b$ are finite monomial constraints with
positive weight one and one ordered factorization. Repeated factors
can either be kept or removed with error $O(\sqrt N)$; they belong to
the bad set for $b<1$.

Here is a direct derivation of the needed mass and its prime-count
specialization. Put $\beta=b/(1+b)$ and fix
$0<\eta<\beta/2$. If $r\le N^{\beta-\eta}$ and $r>q^b$, then
\[
 rq<r^{1+1/b}\le N^{1-(1+1/b)\eta}.
\]
All such bad products therefore cost $o(\mathcal H_N)$ in the prime
count, for each fixed $\eta$. If $r\ge N^\beta$, then
$q\le N/r\le r^{1/b}$; hence the bad condition is automatic on
$r\le q$, $rq\le N$, apart from its equality boundary. The PNT and
partial summation give
\[
 \sum_{\substack{N^\beta<r\le\sqrt N\\r\ {\rm prime}}}
       \bigl(\pi(N/r)-\pi(r)\bigr)
 =\left(\int_\beta^{1/2}\frac{dt}{t(1-t)}+o(1)\right)
      \frac{N}{\log N}.
\]
Indeed the sum of $\pi(r)$ is $O(N/\log^2N)$, while uniformly
$\pi(N/r)=(1+o(1))N/[r\log(N/r)]$ on this fixed exponent interval.
The equality boundary contributes at most $O(\sqrt N)$ pairs.
The intervening strip $N^{\beta-\eta}<r<N^\beta$ is dominated by the
positive domain $r\le q$, $rq\le N$, whose PNT mass is
\[
 \left(\int_{\beta-\eta}^{\beta}\frac{dt}{t(1-t)}+o(1)\right)
 \frac{N}{\log N}=\bigl(O_b(\eta)+o(1)\bigr)\frac{N}{\log N}.
\]
It is a fixed positive-cutoff product family for each fixed $\eta$.
Applying the required positive-product theorem to both the main
region and this outer strip bounds their prime counts by $K$ times
their respective mass coefficients, in units $\mathcal H_N$.
First let $N\to\infty$ with $\eta$ and all sieve parameters fixed,
and then let $\eta\downarrow0$. The resulting coefficient is
\[
 J_b=\int_{b/(1+b)}^{1/2}\frac{dt}{t(1-t)}=\log(1/b).
\]
The ordering $r\le q$ counts each composite just once; squares cost
$O(\sqrt N)$ and require no extra symmetry factor. Thus a proved
positive-product upper theorem with coefficient $K$ gives
$KJ_b\mathcal H_N$, without a new Chen lower-sieve calculation.
Since $\mathcal C(N)$ has a fixed positive
lower bound, the preceding errors and a possible unit are
$o(\mathcal H_N)$.

\paragraph{Transfer and exact arithmetic.}
Each nonunit Chen representation is either in $D_b$ or in $B_b$.
Consequently Li's cited theorem and
\eqref{eq:external-lower-k-required-upper} imply
\begin{equation}\label{eq:external-lower-k-transfer}
 \liminf_{\substack{N\to\infty\\2\mid N}}
 \frac{D_b(N)}{\mathcal H_N}
 \ge \frac{1233}{625}-\frac{67619}{10000}\log(1/b).
\end{equation}
For $0<b<1$, set $u=(1-b)/(1+b)$. A positive series gives the elementary
strict upper bound
\[
 \log(1/b)<2\left(u+\frac{u^3}{3}
                     +\frac{u^5}{5(1-u^2)}\right).
\]
For the requested $b=467/625=0.7472$, this yields
\[
 \log(625/467)<\frac{51919349923661}{178158313856250},
 \qquad
 \frac{1233}{625}-\frac{67619}{10000}\log(625/467)
 >\frac{3972693268066841}{1781583138562500000}>\frac1{500}.
\]
Thus this route gives $D_{0.7472}(N)\ge\mathcal H_N/500$ eventually,
subject to the precise project upper theorem above.
The same rational calculation also permits
\[
 \begin{array}{c|c|c}
 b&\text{strict elementary lower margin}&\text{eventual count}\ \hline
 7471/10000&
 \displaystyle\frac{52783078628162452078269}
 {39841164849867281000000000}&\mathcal H_N/1000\\[4pt]
 747/1000&
 \displaystyle\frac{167157268015739833}
 {398289921308100000000}&\mathcal H_N/2500
 \end{array}
\]
An independent twelve-term positive-series enclosure is recorded in the
companion JSON. Decimal evaluation is only explanatory:
the respective sharper margins are greater than
$0.00223092707$, $0.00132590131$, and $0.00042075440$.
The exact critical expression is
$\exp(-19728/67619)$, approximately $0.7469535199$.
Only exponents strictly above this threshold follow from these two
coefficients; no equality-endpoint assertion is made.
Every eventual bound means that there exists a threshold such that it
holds for every even integer above that threshold. The strict rational
surplus absorbs the $o(1)$ and unit errors; it gives no effective threshold.

\paragraph{Where the original work remains.}
The logical dependency diagram is
\[
 \begin{array}{c}
 \text{Li Theorem 1.1}\ \Downarrow\
 \text{total Chen lower count}
 \end{array}
 \qquad
 \begin{array}{c}
 \text{actual finite upper weights and their factorization}\\
 +\ \text{uniform distribution for their signed atoms}\\
 +\ \text{density certificate and boundary/error control}\ \Downarrow\
 \text{bad-semiprime upper count}
 \end{array}
 \quad\Longrightarrow\quad
 \text{the transfer}.
\]
On the right one must still justify the original compiled replacement
tree, its full and selected prefix gates and pointwise upper property;
the expansion of the actual signed sieve atoms into the coefficient
classes admitted by the cited dispersion theorems; moving residue
$N\pmod d$ uniformly in $N$; the exact coprime main mass, HB grouping,
rough Siegel--Walfisz estimates and Mellin restrictions; positive radial
outer approximation before signed distribution; and the finite
parameter and limit order. These are analytic obligations, not
consequences of positive rational quadrature results.

The numerical part also has its own boundary: the saved fifth-domain
gain must be disjoint from the earlier gain and compatible with the
sixth tree. The recorded replacement arithmetic gives
\[
 K_{\rm end}\le
 \frac{2444205458396424207668787105}{361468820064774123237146624}
 <\frac{67619}{10000}.
\]
This audit reconstructs that subtraction and strict rounding, accepting
the separately replayed density and geometry certificates. It does not
reintegrate them or infer distribution from this fraction. The scope of
the original positive-product theorem is a finite positive sum of
fixed tuple domains at a fixed positive factor cutoff, not arbitrary
bounded arithmetic reweightings, a cutoff vanishing with $N$, or an
infinite adaptive tree.

Once the total lower count is taken directly from Li, this branch does
not use the project's eleven-term Chen lower-weight proof, its
W1/W2 recursions, finite H/h tables, or the 184-generation lower
functional. In particular no table correction may be added to $1.9728$
merely because it improved the separate project lower coefficient.

\paragraph{Meaning of the earlier proof-interface objections.}
The files \texttt{external-chen-input-resolution.tex} and
\texttt{external-chen-input-independent-review.json} establish that
certain separate old-level and new-level inequalities do not by
themselves imply the mixed saving used in the checked derivation.
They also exhibit rational points outside a stated largest-prime
restriction where a printed envelope would affect the integrand.
These are concrete nonimplication or applicability witnesses. They are
not actual prime-count counterexamples to Theorem~1.1, do not prove its
numerical conclusion false, and do not constitute a recomputation of its
full proof with all repairs. This distinction does not erase the
unresolved proof-interface questions.

It is mathematically legitimate to use the final published unconditional
theorem as a named external theorem rather than independently reprove
it. That adds no RH-type conjectural hypothesis. The honest audit
statement is nevertheless: \emph{the final Li theorem is accepted as an
external input, its proof has not been independently closed here, and
the smaller removal coefficient requires the project's separate
analytic theorem}. If that original theorem is established using its
stated unconditional inputs, the resulting implication is an
unconditional corollary with citations. This narrow review certifies
the citation scope, normalization, specialization and exact transfer;
it does not certify the full analytic dependency chain or literature
priority.
